\titre{Rattrapage METNU }
\theme{Méthodes numériques}
\auteur{}
\organisation{AMSCC}

\contenu{



Soit 
\[
A = 
\begin{pmatrix}
5 & 1 & 2 \\
1 & 4 & -1 \\
1 & -1 & 3 
\end{pmatrix}
\quad \text{et} \quad
b = 
\begin{pmatrix}
1 \\
2 \\
3 
\end{pmatrix}.
\]
\begin{enumerate}
    \item Exprimer matriciellement la suite des itérés de la méthode de Jacobi sous la forme \( x^{(k+1)} = F(x^{(k)}) \).
    \item Justifier la convergence de la suite \( (x^{(k)}) \) vers la solution \( x \) de l'équation \( Ax = b \).
    \item Soit \( e^{(k)} = x^{(k)} - x \) l'erreur commise à l'étape \( k \). Donner une matrice \( B \) telle que \( e^{(k)} = B e^{(k-1)} \).
    \item Est-il possible que la matrice \( B \) ait une valeur propre \( \lambda < -1 \) ? Justifier.
\end{enumerate}






















}